\documentclass[11pt]{article}

\usepackage[a4paper,margin=1in]{geometry}
\usepackage{amsmath,amssymb,amsthm}
\usepackage{booktabs}
\usepackage{hyperref}
\usepackage{microtype}

\newtheorem{theorem}{Theorem}
\newtheorem{proposition}[theorem]{Proposition}
\newtheorem{corollary}[theorem]{Corollary}
\newtheorem{lemma}[theorem]{Lemma}
\theoremstyle{remark}
\newtheorem{remark}[theorem]{Remark}

\title{diff-gpt: autoregressive time-series forecasting via \\
       derivative tokenization with MASH $\Sigma\Delta$-quantization}
\author{}
\date{\today}

\begin{document}
\maketitle

\begin{abstract}
We study autoregressive forecasting of continuous time series with a
decoder-only transformer trained on a derivative-based tokenization.
The encoder differentiates the input to order $k \ge 1$, scales by its
observed range, and quantizes with a \emph{MASH} (multi-stage
noise-shaping) $k$-th-order $\Sigma\Delta$ modulator; the decoder
inverts this. We prove two tokenization-level bounds, both independent
of the model: the achievable cross-entropy loss grows linearly in
$\log V$, and the reconstruction error is uniformly bounded by $\Delta/2
= D/(V-2-2^k)$ for any $k$ such that $V > 2^k + 2$. We further observe
that while MASH's reconstruction bound is tight against pure
quantization noise, \emph{model prediction error} compounds as
$O(\Delta \cdot T^{k-1})$ through $k$ decoder integrations --- giving a
principled reason to prefer $k = 1$ for forecasting despite the clean
higher-order theory. We introduce an ordinal Gaussian soft cross-entropy
that converts bin-wise classification into a distance-aware loss over
ordinal bins; on M4 Hourly it reduces sMAPE from $18.58$ to $17.49$
(6\% relative) without architectural change.
\end{abstract}

\section{Introduction}

Most state-of-the-art time-series forecasters are regression models with
specialized inductive biases (patching, inverted attention, channel
mixing) trained against an $\ell_2$ or Huber loss. By contrast, large
language models are decoder-only transformers trained with cross-entropy
over a discrete vocabulary. The question we address is whether a pure
language-model recipe --- fixed architecture, fixed loss,
autoregressive generation --- is competitive on standard forecasting
benchmarks, given a faithful discrete representation of the continuous
input.

The key choice is the tokenization. We describe a derivative-based
encoder with MASH $\Sigma\Delta$ quantization that satisfies two closed-
form vocabulary-scaling laws: (a) the ideal cross-entropy loss grows
linearly in $\log V$ (Prop.~\ref{prop:ce}) and (b) the reconstruction
error between a round-trip of encoder and decoder is bounded uniformly
in $T$ by $\Delta/2$ for any order $k$ (Thm.~\ref{thm:recon}). Both are
tokenization-level statements --- they hold regardless of what model
sits on top.

The rest of the paper. Section~\ref{sec:encoder} describes the encoder
and decoder; Section~\ref{sec:theory} proves the two scaling laws and
warns that the reconstruction theorem's tight bound is against
quantization noise only --- the model's prediction error compounds
less kindly at higher $k$; Section~\ref{sec:loss} introduces ordinal
soft cross-entropy and motivates it from the differential-entropy
bound; Section~\ref{sec:experiments} reports M4 Hourly and ETTh1 results.

\section{Derivative tokenization}
\label{sec:encoder}

Let $x \in \mathbb{R}^{T \times F}$ be a multivariate signal with $F$
channels of length $T$. We fix a derivative order $k \ge 1$ and define
$D_c = \max_{t} |x_{t+1,c} - x_{t,c}|$ per channel (the \emph{domain
of definition}). For even $V$ with $V > 2^k + 2$, define the per-channel
scale
\begin{equation}
s_c \;=\; \frac{V - 2 - 2^k}{2 D_c}, \qquad \Delta_c \;=\; \frac{1}{s_c} \;=\; \frac{2 D_c}{V - 2 - 2^k}.
\end{equation}
The $-2^k$ headroom is the maximum overshoot MASH can produce over
the unshaped input (Lemma~\ref{lem:mash-bound}); without it the
quantizer would saturate at signal peaks and lose the noise-shaping
guarantee.

\subsection{MASH $k$-th-order modulator.}
Each of the $k$ stages is a stable first-order $\Sigma\Delta$ loop.
Stage 1 operates on the scaled $k$-th finite difference $y_t = s \cdot
\Delta^k x_t$; stage $i \ge 2$ operates on $-\varepsilon^{(i-1)}_t$, the
previous stage's quantization error (bounded by $\tfrac{1}{2}$ in bin
units). Let $q^{(i)}_t$ denote stage $i$'s integer output and
$\varepsilon^{(i)}_t = q^{(i)}_t - (\text{stage $i$'s quantizer input})$.

The final token stream is the digital-differentiator combination
\begin{equation}
q_t \;=\; \sum_{i=0}^{k-1} (1 - z^{-1})^i \, q^{(i+1)}_t \;=\;
           q^{(1)}_t + \Delta q^{(2)}_t + \Delta^2 q^{(3)}_t + \cdots
           + \Delta^{k-1} q^{(k)}_t,
\label{eq:mash-combined}
\end{equation}
shifted by $V/2$ and clipped to $\{0, \dots, V-1\}$.

\subsection{Decoder.}
Given the start prefix $(x_1, \ldots, x_k)$ and token sequence $(q_t)$,
the decoder reconstructs the signal by dequantizing and integrating
$k$ times:
\begin{equation}
\hat x_{t+1} \;=\; \hat x_t + (q_t - V/2)\,\Delta.
\end{equation}

\subsection{Per-column orders.}
For multivariate input where channels have qualitatively different
dynamics (e.g., a slowly-drifting trend column alongside a
stationary-noise column), a single shared $k$ is suboptimal. The
implementation accepts $k$ as a per-column vector $k_c$ instead of a
scalar. Each column is differentiated to its own $k_c$, scaled by
$(V - 2 - 2^{k_c}) / (2 D_c)$ if $k_c \ge 2$ (for MASH headroom) or
$(V - 2) / (2 D_c)$ otherwise, then passed through the appropriate
$\Sigma\Delta$ variant per column (plain round for $k_c = 0$, carry
scheme for $k_c = 1$, MASH for $k_c \ge 2$). The encoder aligns all columns to
the shared valid window $T - \max_c k_c$, stores the raw prefix as
$start$, and the decoder reconstructs each column using only its own
$k_c$ cumsums. Theorem~\ref{thm:recon} and Prop.~\ref{prop:ce} apply
per column with $D_c$, $k_c$ replacing their scalar counterparts.

\section{Tokenization contract}
\label{sec:theory}

Two tokenization-level results, both independent of the model.

\subsection{Reconstruction bound}

\begin{theorem}[MASH reconstruction, any $k$]
\label{thm:recon}
For any signal with $D = \max_t |\Delta^k x_t| > 0$ and vocabulary size
$V > 2^k + 2$, the MASH encoder--decoder pair satisfies
\begin{equation}
\|\hat x - x\|_\infty \;\le\; \frac{D}{V - 2 - 2^k} \;=\; \frac{\Delta}{2},
\qquad \text{uniformly in } T.
\end{equation}
\end{theorem}

The proof relies on a standard identity about the cascade.

\begin{lemma}[MASH noise-shaping]
\label{lem:mash-bound}
Let $\varepsilon_k := \varepsilon^{(k)}$ be the quantization error of the
final stage. The combined output \eqref{eq:mash-combined} satisfies
\begin{equation}
q_t - y_t \;=\; (1 - z^{-1})^k \, \varepsilon_{k,t}, \qquad
|\varepsilon_{k,t}| \;\le\; \tfrac{1}{2}.
\end{equation}
\end{lemma}

\begin{proof}[Sketch]
Induction on the cascade. Stage 1's first-order $\Sigma\Delta$ gives
$q^{(1)}_t = y_t + \Delta \varepsilon^{(1)}_t$. Stage $i \ge 2$ with
input $-\varepsilon^{(i-1)}_t$ gives $q^{(i)}_t = -\varepsilon^{(i-1)}_t
+ \Delta \varepsilon^{(i)}_t$. Substituting into \eqref{eq:mash-combined}
and telescoping the intermediate $\varepsilon^{(i)}$ terms yields
$q_t = y_t + \Delta^k \varepsilon_{k,t}$. The bound
$|\varepsilon_{k,t}| \le 1/2$ is immediate: each 1st-order stage rounds
to the nearest integer, so its internal quantization error is in
$(-\tfrac{1}{2}, \tfrac{1}{2}]$.
\end{proof}

\begin{lemma}[bounded MASH output]
\label{lem:mash-range}
$|q_t - y_t| \le 2^{k-1}$.
\end{lemma}

\begin{proof}
$|\Delta^k \varepsilon_k| \le \|(1-z^{-1})^k\|_1 \cdot \max|\varepsilon_k|
= 2^k \cdot \tfrac{1}{2} = 2^{k-1}$.
\end{proof}

The $-2^k$ headroom in the scale's denominator ensures this overshoot
stays within the quantizer range, avoiding saturation.

\begin{proof}[Proof of Theorem~\ref{thm:recon}]
The decoded increment is $(q_t - V/2)\Delta = (y_t + \Delta^k
\varepsilon_{k,t})\Delta$. The true increment is $y_t \Delta$. Their
difference is $\Delta^k \varepsilon_{k,t} \cdot \Delta$. Summing $k$
times from the exact starts telescopes the $\Delta^k$ down to
$\varepsilon_{k,t}$:
\begin{equation}
\hat x_T - x_T \;=\; \varepsilon_{k,T-k} \cdot \Delta.
\end{equation}
By Lemma~\ref{lem:mash-bound}, $|\varepsilon_k| \le \tfrac{1}{2}$, so
$|\hat x_T - x_T| \le \Delta/2$.
\end{proof}

\begin{corollary}[scaling]
$\|\hat x - x\|_\infty = O(D/V)$. Vocabulary doubling halves the
reconstruction error, uniformly in $T$ and for any $k$.
\end{corollary}

\subsection{Prediction-error compounding vs. quantization noise}
\label{sec:pred-compound}

Theorem~\ref{thm:recon} bounds the residual between $\hat x$ and $x$
\emph{given} a perfect encoder/decoder round-trip. In autoregressive
forecasting, the decoder does not consume the exact encoded tokens of
the true signal; it consumes tokens \emph{predicted} by the model,
which carry prediction error on top of quantization noise.

\begin{proposition}[prediction-error amplification at order $k$]
\label{prop:pred}
Suppose the model's per-step prediction error has typical magnitude
$\eta$ bins (so $\eta \cdot \Delta$ in value space per step). After $k$
decoder integrations the forecast error at the horizon end grows as
\begin{equation}
|\hat x_{T+H} - x_{T+H}| \;=\; O\!\left(\eta \cdot \Delta \cdot H^{\,k-1}\right),
\end{equation}
where $H$ is the number of forecast steps.
\end{proposition}

\begin{proof}[Sketch]
Token prediction error $\delta q_t$ of typical size $\eta$ enters the
decoder as diff-estimate error $\delta q_t \cdot \Delta$. Under $k$
successive cumsum operations, an i.i.d.\ or coherent bias sequence of
size $\eta \Delta$ produces output whose magnitude at step $H$ scales
as $\eta \Delta H^{k-1}$ (integration of polynomial in $H$).
\end{proof}

\begin{remark}[$k=1$ preferred for forecasting]
At $k=1$ the amplification is $H^0 = 1$: prediction errors contribute
proportionally to their bin size, not multiplicatively across the
horizon. At $k=2$ a single-bin token error grows $H \cdot \Delta$ over
a 48-step forecast, on the order of $24 \Delta$ --- much larger than
the $\Delta/2$ quantization bound. Empirically (Section~\ref{sec:experiments})
this manifests as catastrophic forecasting regression at $k=2$ despite
strictly lower token-level cross-entropy.
\end{remark}

\subsection{Cross-entropy bound}

Let $X_t$ be the continuous next-step increment conditional on context
$\mathrm{ctx}$, with density $p(\cdot \mid \mathrm{ctx})$ and differential
entropy $h(X_t \mid \mathrm{ctx}) = -\int p \log p$.

\begin{proposition}[achievable CE]
\label{prop:ce}
In the fine-quantization limit $\Delta \to 0$, the expected cross-
entropy of the optimal autoregressive token predictor satisfies
\begin{equation}
\mathcal{L}^\star(V) \;=\; h(X_t \mid \mathrm{ctx}) + \log V - \log(2D) \;+\; o(1).
\end{equation}
Equivalently, $\mathcal L^\star$ is linear in $\log V$ with slope $1$.
\end{proposition}

\begin{proof}[Sketch]
Standard quantized-entropy relation (Cover \& Thomas, \S 8.3): for a
continuous RV with density $p$ and bin width $\Delta$,
\begin{equation}
H(\mathrm{bin}) \;=\; h(X) - \log \Delta + o(1), \qquad \Delta \to 0.
\end{equation}
Condition on $\mathrm{ctx}$ and substitute $\Delta = 2D/(V - 2 - 2^k)
\sim 2D/V$.
\end{proof}

\begin{corollary}[vocabulary-invariant skill]
\label{cor:skill}
Two models at vocabulary sizes $V_1 \neq V_2$ are comparable only
through the \emph{differential-entropy residual}
\begin{equation}
\tilde h(V) \;:=\; \mathcal L(V) - \log V + \log(2D),
\end{equation}
which estimates $h(X_t \mid \mathrm{ctx})$ and is independent of $V$ at
the rate.
\end{corollary}

\subsection{Empirical check}

M4 Hourly, global model, $V$-sweep with $\sigma$ scaled $\propto V$ to
hold value-space soft-target smoothing constant:

\begin{center}
\begin{tabular}{lcccc}
\toprule
$V$ & $\sigma$ & val CE (nats) & $\Delta \log V$ & $\Delta \text{(CE)}$ \\
\midrule
64   & 0.5 & 1.01 & --- & --- \\
128  & 1.0 & 1.70 & $\log 2 = 0.69$ & $0.69$ \\
256  & 2.0 & 2.35 & $\log 2 = 0.69$ & $0.65$ \\
512  & 4.0 & 3.04 & $\log 2 = 0.69$ & $0.69$ \\
\bottomrule
\end{tabular}
\end{center}

Empirical CE increments match $\log 2$ per vocab doubling to within
$0.04$ nat. Consistent with Prop.~\ref{prop:ce}.

\section{Ordinal soft cross-entropy}
\label{sec:loss}

Standard cross-entropy treats the $V$ bins as nominal categories: a
prediction off by one bin is penalized identically to one off by
fifty. Bins are an ordinal index over a continuous scale --- off by one
is almost right.

We replace the one-hot target $\mathbf 1_{i = y}$ with a discrete
Gaussian over bin indices:
\begin{equation}
p^{\text{soft}}_i(y; \sigma) \;\propto\; \exp\!\left(-\tfrac{(i - y)^2}{2\sigma^2}\right),
\qquad \sigma \ge 0,
\end{equation}
and train with soft-target cross-entropy $\mathcal L_\sigma = -\sum_i
p^{\text{soft}}_i \log q_i$, where $q = \mathrm{softmax}(\text{logits})$.
Taking $\sigma \to 0$ recovers ordinary cross-entropy with one-hot
targets.

\paragraph{Why this matters for a $\Sigma\Delta$-quantized signal.}
A continuous value landing near a bin boundary is, modulo the MASH
carry, assigned to one of two adjacent bins; the encoder can shift
either way depending on the running accumulator. One-hot training
pretends this assignment is deterministic; soft targets acknowledge
the effectively stochastic tie-breaking. More broadly, any aleatoric
component of $p(X \mid \mathrm{ctx})$ makes the true conditional
distribution over bins a smooth Gaussian-like shape; label smoothing
with scale $\sigma$ in bin units is the ordinal equivalent.

\paragraph{Bin-space versus value-space.}
Because bin width $\Delta \approx 2D/V$, a constant value-space
smoothing $\sigma_{\text{val}} = \sigma \cdot \Delta$ requires $\sigma$
to scale linearly in $V$. A value of $\sigma$ tuned at one $V$ is
under-scaled at larger $V$ by the same factor.

\section{Experiments}
\label{sec:experiments}

\subsection{M4 Hourly, short-term global protocol}

Following the M4 leaderboard: per-series $z$-score normalization, one
global model across 414 Hourly series, horizon $H = 48$, best-val
checkpoint restored. Model: 2-layer GPT, $d_{\text{model}} = 64$,
$h = 4$, block 138, 20k iters, AdamW, temperature-0 argmax at inference.

\paragraph{Soft-CE sweep at $V = 256$.}
\begin{center}
\begin{tabular}{lcc}
\toprule
$\sigma$ & sMAPE & MASE \\
\midrule
$0$ (plain CE) & 18.58 & --- \\
$1.0$ & 17.55 & 1.663 \\
$2.0$ & \textbf{17.49} & \textbf{1.659} \\
$3.0$ & 17.35 & 1.949 \\
\bottomrule
\end{tabular}
\end{center}

$\sigma = 2$ is the joint-metric optimum. At $\sigma = 3$ sMAPE
continues to improve but MASE degrades sharply --- a classic
over-smoothing signature where forecasts regress toward the
conditional mean, rewarded by symmetric sMAPE but penalized by the
naive-seasonal-normalized MASE.

\paragraph{Vocabulary sweep (U-curve).}
\begin{center}
\begin{tabular}{cccc}
\toprule
$V$ & $\sigma$ & sMAPE & MASE \\
\midrule
64  & 0.5 & 19.96 & 2.019 \\
128 & 1.0 & 19.39 & 1.888 \\
\textbf{256} & \textbf{2.0} & \textbf{17.49} & \textbf{1.659} \\
512 & 4.0 & 18.21 & 1.772 \\
\bottomrule
\end{tabular}
\end{center}

Non-monotone with interior optimum at $V = 256$. Smaller $V$ hurts
reconstruction precision (bin width doubles); larger $V$ over-refines
for the $\sim 280$k available training tokens. The CE bound is always
tight ($\log 2$ per doubling) but the prediction-quality curve is not.

\paragraph{Order-of-derivative ablation.}
\begin{center}
\begin{tabular}{ccccc}
\toprule
$k$ & val CE & sMAPE & MASE & note \\
\midrule
\textbf{1} & \textbf{2.35} & \textbf{17.49} & \textbf{1.66} & --- \\
2 & 2.28 & 39.27 & 8.93 & prediction-error compounding \\
\bottomrule
\end{tabular}
\end{center}

Token-level training loss is \emph{lower} at $k=2$, confirming that
MASH's higher-order tokenization is learnable. But reconstruction via
double integration amplifies model prediction error linearly with the
forecast horizon (Prop.~\ref{prop:pred}), collapsing sMAPE.
Practically, $k = 1$ is preferred for autoregressive forecasting.

\paragraph{Model-size sensitivity ($V = 256$, $\sigma = 2$).}
\begin{center}
\begin{tabular}{cccccc}
\toprule
$d_{\text{model}}$ & $n_{\text{layer}}$ & params & val CE & sMAPE & MASE \\
\midrule
32 & 2 & $\sim$30k  & 2.36 & 18.47 & 1.88 \\
\textbf{64} & \textbf{2} & $\sim$130k & 2.35 & \textbf{17.49} & \textbf{1.66} \\
64 & 4 & $\sim$260k & 2.35 & 17.41 & 1.68 \\
128 & 2 & $\sim$400k & 2.38 & 18.67 & 1.99 \\
\bottomrule
\end{tabular}
\end{center}

With $\sim 280$k training tokens, the $64 \times 2$ setting
($\sim 2$ tokens/parameter) is near the data-limited ceiling. Wider
(128) overfits (best val at step 2000 of 20000); deeper (64$\times$4)
matches within noise at $2.3\times$ training cost.

\subsection{ETTh1, long-term multivariate protocol}

iTransformer protocol: 7 channels, seq\_len = 96, pred\_len = 96,
chronological 12/4/4-month split, StandardScaler from train. Single
model, $d = 128$, $h = 4$, $n_{\text{layer}} = 2$, 2000 iters, batch 32.

\begin{center}
\begin{tabular}{lcc}
\toprule
config & MSE & MAE \\
\midrule
$V = 256$, plain CE                 & 0.9861 & 0.6056 \\
$V = 64$, plain CE                  & 0.9787 & 0.6030 \\
$V = 64$, soft-CE $\sigma = 1.0$    & 0.9806 & 0.6034 \\
\bottomrule
\end{tabular}
\end{center}

ETTh1 is dominated by a capacity/data-fit floor around MSE $= 0.98$:
vocab and soft-CE move the number by within-noise amounts. Contrast
with M4 suggests the \emph{value} of soft targets concentrates in
data-starved regimes where many short series force extrapolation of
conditional distributions without enough signal to localize a single
bin.

\subsection{Vocabulary / loss scaling, summary}
\begin{center}
\begin{tabular}{lcc}
\toprule
                         & per vocab doubling & reference \\
\midrule
ideal CE $\mathcal L^\star(V)$   & $+\log 2 \approx 0.69$ nats & Prop.~\ref{prop:ce} \\
reconstruction $\|\hat x - x\|_\infty$ & $\times \tfrac12$ & Thm.~\ref{thm:recon} \\
ordinal smoothing $\sigma^\star$ & $\times 2$ (hold value-space constant) & \S\ref{sec:loss} \\
\bottomrule
\end{tabular}
\end{center}

\section{Related work}

\textbf{Inverted/channel-mixing transformers.} iTransformer
\cite{itransformer2024} and PatchTST \cite{patchtst2023} remove the
autoregressive head and regress the entire horizon in one shot, the
dominant design on long-horizon benchmarks. Our model keeps the GPT
recipe unchanged; the tokenization carries the inductive bias.

\textbf{Ordinal and label-smoothed classification.} Label smoothing
\cite{szegedy2016} distributes probability uniformly over non-target
classes. Distance-weighted variants (cumulative-link ordinal softmax)
appear in the ordinal-regression literature; Gaussian soft targets
specifically in CORAL \cite{cao2020} and in diffusion-model score
matching of discretized targets \cite{austin2021}. Our contribution is
situating the smoothing scale $\sigma$ in a $\Sigma\Delta$-tokenized
vocabulary.

\textbf{$\Sigma\Delta$ quantization.} First-order sigma--delta
modulators are classical in analog-to-digital conversion
\cite{candy1992}; MASH (multi-stage noise shaping) extends to arbitrary
order with stability guarantees, originally for audio/oversampling
applications \cite{hayashi1986}. We apply MASH as a differentiable-free
preprocessing stage in front of a language model.

\section{Future research directions}
\label{sec:future}

The experiments here focus on deterministic sMAPE/MSE benchmarks.
diff-gpt's architectural choices --- autoregressive language-model
recipe over a quantized vocabulary --- give it structural advantages in
a different regime, which we believe deserves dedicated follow-up.

\paragraph{Probabilistic forecasting (quantile / interval prediction).}
The decoder emits a full distribution over next-token bins at every
step; $N$-sample rollouts give native P10/P50/P90 estimates without
a pinball-loss head or quantile-regression machinery. Direct
competitors (iTransformer, N-BEATS) optimize MSE, i.e.\ the conditional
mean. Applications: energy-load dispatch with confidence bands,
financial risk (VaR, expected shortfall), reliability engineering,
weather ensembles, inventory planning.

\paragraph{Heavy-tailed and multi-modal conditional distributions.}
$\ell_2$ regression collapses toward the conditional mean and fails
catastrophically on bursty / spike / regime-change series where the
true $p(X_t \mid \mathrm{ctx})$ is bimodal (``jump up or jump down,
not in between''). A categorical head models this directly. Candidate
domains: network traffic, high-frequency finance (CDS spreads,
order-book imbalances), industrial sensors with fault modes.

\paragraph{Foundation-style zero-/few-shot forecasting.}
The LM recipe admits trivial pretrain-then-transfer: one model
pretrained on a million heterogeneous series, few-shot on a new one.
This is the angle pursued by Chronos, TimesFM, MOIRAI. Our
tokenization contract (two closed-form scaling laws) is theoretically
cleaner than their binning schemes, which is a plausible advantage at
scale.

\paragraph{Heterogeneous / multi-resolution series in one model.}
The encoder self-normalizes per channel via $D$, letting a single
parameter set handle $\mu V$ ECG, $m/s^2$ seismograms, and dollar-scale
macro in one training run. Regression baselines need per-channel heads
or manual unification. For foundation models this is a core problem
and our scheme is structurally simpler.

\paragraph{Anomaly detection via conditional likelihood.}
$-\log p(\text{token}_t \mid \mathrm{ctx})$ is available for free at
each step and calibrates correctly under variable conditional
variance. Regression-based approaches rely on $\|y_t - \hat y_t\|$,
which miscalibrates on heteroscedastic series. Candidate domains:
industrial monitoring, network intrusion, patient monitoring.

\paragraph{Conditional / counterfactual generation.}
``What if $X_t$ jumps by $\Delta$ next week?'' reduces to token
substitution in the prefix, then resume generation. Regression models
need a separate conditional mode or structured covariate mechanism.

\paragraph{Edge inference.}
A 130k-parameter model attains the results in Section~\ref{sec:experiments};
the bulk of the inductive bias lives in the tokenization, not the
weights. This makes sub-mobile-CPU inference realistic for production
forecasting on IoT or embedded devices --- a regime where Transformer
SOTA typically does not fit.

\paragraph{Where the architecture will not win.}
Deterministic MSE benchmarks at long horizons (ETTh1 / Traffic at
pred\_len $\ge 336$) are structurally unfavorable: one-shot inverted
transformers emit the whole horizon without autoregressive error
compounding, and Prop.~\ref{prop:pred} guarantees our error grows at
rate $H^{k-1}$. Deterministic long-horizon forecasting is not a
target; the wins above are.

\section{Conclusion}

A decoder-only transformer trained autoregressively on derivative-plus-
MASH-quantized tokens gives competitive short- and long-term forecasts
without a specialized head. The tokenization contract provides two
closed-form vocabulary-scaling laws valid at any derivative order;
prediction-error compounding gives a principled preference for $k=1$
in forecasting despite MASH's tight higher-order reconstruction
bound. Ordinal soft cross-entropy matches the natural bin-wise geometry
and yields a 6\% relative sMAPE reduction on M4 Hourly without
architectural change.

\section*{Code availability}

Source, benchmarks, and tests are at
\url{https://github.com/fxeqxmulfx/diff-gpt}. Install via
\texttt{pip install git+https://github.com/fxeqxmulfx/diff-gpt}.

\begin{thebibliography}{9}
\bibitem{itransformer2024} Liu, Y.\ et al.\ \emph{iTransformer: Inverted
Transformers are Effective for Time Series Forecasting}. ICLR 2024.
arXiv:2310.06625.

\bibitem{patchtst2023} Nie, Y.\ et al.\ \emph{A Time Series is Worth 64
Words: Long-term Forecasting with Transformers}. ICLR 2023.

\bibitem{szegedy2016} Szegedy, C.\ et al.\ \emph{Rethinking the
Inception Architecture for Computer Vision}. CVPR 2016.

\bibitem{cao2020} Cao, W., Mirjalili, V., Raschka, S.\ \emph{Rank-
consistent Ordinal Regression for Neural Networks with Application to
Age Estimation}. Pattern Recognition Letters 2020.

\bibitem{austin2021} Austin, J.\ et al.\ \emph{Structured Denoising
Diffusion Models in Discrete State-Spaces}. NeurIPS 2021.

\bibitem{candy1992} Candy, J.C., Temes, G.C.\ \emph{Oversampling Delta-
Sigma Data Converters}. IEEE Press 1992.

\bibitem{hayashi1986} Hayashi, T., Kaneko, Y., Morimoto, A.,
Sugawara, T., Iwasaki, K.\ \emph{A multistage delta-sigma modulator
without double integration loop}. IEEE ISSCC 1986.
\end{thebibliography}

\end{document}
